% ============================================================
% COLORS, CHAPTER PALETTE AND CHAPTER STRUCTURE
% ============================================================

% Palette sortiert nach max. 3-Fenster-Kontrast (CIEDE2000)
% ChapterColor0 bleibt ETH-Rot für den Frontchapter
\definecolor{ChapterColor0}{HTML}{A50D0D}
\definecolor{ChapterColor0Light}{HTML}{EED2D2}
\definecolor{ChapterColor1}{HTML}{5B8C00}
\definecolor{ChapterColor1Light}{HTML}{DDE8B5}
\definecolor{ChapterColor2}{HTML}{7211D4}
\definecolor{ChapterColor2Light}{HTML}{E0D2EE}
\definecolor{ChapterColor3}{HTML}{D44B11}
\definecolor{ChapterColor3Light}{HTML}{EEDAD2}
\definecolor{ChapterColor4}{HTML}{00897B}
\definecolor{ChapterColor4Light}{HTML}{B2DFDB}
\definecolor{ChapterColor5}{HTML}{2B0DA5}
\definecolor{ChapterColor5Light}{HTML}{D8D2EE}
\definecolor{ChapterColor6}{HTML}{A5680D}
\definecolor{ChapterColor6Light}{HTML}{EEE3D2}
\definecolor{ChapterColor7}{HTML}{2E7D32}
\definecolor{ChapterColor7Light}{HTML}{C8E6C9}
\definecolor{ChapterColor8}{HTML}{D411C0}
\definecolor{ChapterColor8Light}{HTML}{EED2EB}
\definecolor{ChapterColor9}{HTML}{1199D4}
\definecolor{ChapterColor9Light}{HTML}{D2E5EE}
\definecolor{ChapterColor10}{HTML}{2BA50D}
\definecolor{ChapterColor10Light}{HTML}{D8EED2}
\definecolor{ChapterColor11}{HTML}{A50D68}
\definecolor{ChapterColor11Light}{HTML}{EED2E3}
\definecolor{ChapterColor12}{HTML}{860DA5}
\definecolor{ChapterColor12Light}{HTML}{E8D2EE}
\definecolor{ChapterColor13}{HTML}{0DA5A5}
\definecolor{ChapterColor13Light}{HTML}{D2EEEE}
\definecolor{ChapterColor14}{HTML}{0D4AA5}
\definecolor{ChapterColor14Light}{HTML}{D2DDEE}
\definecolor{ChapterColor15}{HTML}{D4114B}
\definecolor{ChapterColor15Light}{HTML}{EED2DA}
\definecolor{ChapterColor16}{HTML}{0DA54A}
\definecolor{ChapterColor16Light}{HTML}{D2EEDD}
\definecolor{ChapterColor17}{HTML}{86A50D}
\definecolor{ChapterColor17Light}{HTML}{E8EED2}
\definecolor{ChapterColor18}{HTML}{1124D4}
\definecolor{ChapterColor18Light}{HTML}{D2D5EE}
\definecolor{ChapterColor19}{HTML}{D4C011}
\definecolor{ChapterColor19Light}{HTML}{EEEBD2}
\definecolor{ETHRot}{HTML}{CC0000}

% Semantische Farben für Text-Box-Bindungen
\definecolor{TextVorBoxColor}{HTML}{0A0A1E}   % Sehr dunkles Blau (einleitend)
\definecolor{TextNachBoxColor}{HTML}{0A1A0A}   % Sehr dunkles Grün (ergänzend)

\newcommand{\chaptercolor}{ChapterColor0}
\newcommand{\chaptercolorlight}{ChapterColor0Light}
\newcommand{\ChapterBarColor}{\chaptercolor}
\newcommand{\SubsectionBarColor}{\chaptercolor!75}
\newcommand{\FormulaboxBackColor}{\chaptercolorlight!40}

% Statische Farben für Box-Titel und neutrale Sub-Boxen (unabhängig vom Kapitel)
\definecolor{StaticBoxTitleBack}{HTML}{E2E8F0} % Helles blaugrau (Tailwind Slate-200)
\definecolor{StaticBoxTitleText}{HTML}{0F172A} % Dunkles blaugrau (Tailwind Slate-900)
\definecolor{StaticSubBoxBack}{HTML}{F8FAFC} % Sehr helles grau (Tailwind Slate-50)
\definecolor{StaticSubBoxBorder}{HTML}{CBD5E1} % Mittleres grau (Tailwind Slate-300)

% Dynamische Anpassung an Kapitel, damit Boxen nicht alle gleich aussehen
\newcommand{\BoxTitleColor}{\chaptercolorlight}
\newcommand{\BoxTitleTextColor}{TextVorBoxColor}
\newcommand{\BoxBackColor}{\chaptercolorlight!8}
\newcommand{\BoxFrameColor}{\chaptercolorlight!25}

\newcommand{\GoalConditionBackColor}{StaticSubBoxBack}
\newcommand{\GoalConditionFrameColor}{StaticSubBoxBorder}
\newcommand{\GoalTargetBackColor}{StaticBoxTitleBack}
\newcommand{\GoalTargetFrameColor}{StaticSubBoxBorder}

\newcount\ZSFchapterPaletteIndex
\newcommand{\ZSFchapterPaletteSize}{20}
\newcommand{\ZSFSetChapterPaletteByIndex}[1]{%
  \renewcommand{\chaptercolor}{ChapterColor#1}%
  \renewcommand{\chaptercolorlight}{ChapterColor#1Light}%
}

\newcommand{\SetChapterPaletteByNumber}{%
  \ZSFchapterPaletteIndex=\value{section}%
  \loop
  \ifnum\ZSFchapterPaletteIndex>\numexpr\ZSFchapterPaletteSize-1\relax
    \advance\ZSFchapterPaletteIndex by -\ZSFchapterPaletteSize
  \repeat
  \ifnum\value{section}>\numexpr\ZSFchapterPaletteSize-1\relax
    \PackageWarning{ZSF}{Chapter palette wrapped at section \thesection. Increase palette if unique colors are required}%
  \fi
  \ZSFSetChapterPaletteByIndex{\the\ZSFchapterPaletteIndex}%
}

\newif\ifzsfAfterChapterBar
\zsfAfterChapterBarfalse
\newcommand{\ZSFMarkAfterChapterBar}{\global\zsfAfterChapterBartrue}
% Kollabiert den before-skip der SubsectionBar auf \ZSFsubsectionAfterChapterGap,
% wenn sie direkt nach einer ChapterBar steht. Reset nach Verwendung.
\newcommand{\ZSFCollapseAfterChapterBar}{%
  \ifzsfAfterChapterBar
    \vspace*{\dimexpr\ZSFsubsectionAfterChapterGap-\ZSFsubsectionBarBeforeSkip\relax}%
    \global\zsfAfterChapterBarfalse
  \fi
}

\makeatletter
\newcommand{\ZSFPostChapterSpacing}{\vspace*{\ZSFpostChapterNudge}}
\newcommand{\StartChapter}[2][]{%
  \ifnum\value{zsfFirstChap}=1\setcounter{zsfFirstChap}{0}\fi
  \refstepcounter{section}
  \SetChapterPaletteByNumber
  \setcounter{subsection}{0}
  \if\relax\detokenize{#1}\relax\else\label{#1}\fi
  \ChapterBar{\thesection. #2}
  \ZSFMarkAfterChapterBar
  \ZSFPostChapterSpacing
}
\makeatother

% Variante: Kapitel auf neuer Spalte beginnen (ersetzt rohes \columnbreak)
\newcommand{\StartChapterOnNewColumn}[2][]{%
  \columnbreak
  \StartChapter[#1]{#2}%
}

\newcommand{\StartFrontChapter}[1]{
  \ifnum\value{zsfFirstChap}=1\setcounter{zsfFirstChap}{0}\else\columnbreak\fi
  \renewcommand{\chaptercolor}{ChapterColor0}
  \renewcommand{\chaptercolorlight}{ChapterColor0Light}
  \setcounter{subsection}{0}
  \ChapterBar{#1}
  \ZSFMarkAfterChapterBar
  \ZSFPostChapterSpacing
}
